% Source: https://artofproblemsolving.com/community/c3h1981823_tanker_truck
% Stable identity: aops-topic-1981823-q01
A tanker truck full of water begins to leak. The truck's fuel consumption per kilometre is proportional to the total weight it carries (the truck itself plus the water), the water leaks out at a constant rate and the truck moves at constant speed; ignore the weight of the fuel. After travelling $200$ km the truck is half full of water and has used half a tank of fuel. Empty of water, the truck would use $\dfrac{1}{6}$ of a tank for the same $200$ km. What fraction of a tank would the truck use for the $200$ km if there were no leak?
